% !TeX root = ../../Book5.tex
% 纲要标题：### 2.3 群代数的分解
% 纲要位置：工作记录/计划纲要.md，第 3775 行。
% 纲要细目（写作范围）：
% * 正则特征标与不可约表示维数的平方和
% * 复群代数的矩阵块分解
% * 中心幂等元及块的半单情形


\section{群代数的分解}
\label{b5:sec:0203}

正则表示包含全部不可约类型，并且它的迹很容易计算。两者结合，把不可约次数与群阶联系起来。随后再用中心幂等元从群代数中直接取出各个等型分量，说明特征标不仅记录数字，也能构造投影。


\subsection{正则特征标}
\label{b5:subsec:0203:01}

\begin{theorem}\label{b5:thm:regular-decomposition}
设 \(G\) 为有限群，\(\chi_1,\ldots,\chi_r\) 为全部不可约复特征标，次数为 \(d_i=\chi_i(1)\)，相应表示为 \(S_i\)。正则特征标满足
\[
\chi_{\rm reg}(g)=
\begin{cases}|G|,&g=1,\\0,&g\ne1,\end{cases}
\qquad
\mathbb C[G]\simeq\bigoplus_iS_i^{\oplus d_i}.
\]
特别地，\(\sum_i d_i^2=|G|\)。
\end{theorem}
\begin{proof}
左平移 \(x\mapsto gx\) 只有在 \(g=1\) 时有不动点，所以置换特征标给出第一式。由\cref{b5:cor:character-irreducibility}，\(S_i\) 在正则表示中的重数为
\[
\langle\chi_i,\chi_{\rm reg}\rangle_G
=\overline{\chi_i(1)}=d_i.
\]
于是得到分解。比较两侧维数即得平方和。
\end{proof}

对每个 \(S_i\) 选定正交归一基，将其表示矩阵写成酉矩阵。由\cref{b5:thm:matrix-coefficient-orthogonality}，函数 \(g\mapsto\rho_i(g)_{ab}\) 两两正交，每个范数平方为 \(1/d_i\)。刚证的平方和说明它们共有 \(|G|\) 个，而全函数空间 \(\mathbb C^G\) 的维数也是 \(|G|\)。因此
\[
\bigl\{\,\sqrt{d_i}\,\rho_i(\,\cdot\,)_{ab}
\mid 1\le i\le r,\ 1\le a,b\le d_i\,\bigr\}
\]
构成 \(\mathbb C^G\) 的标准正交基，每个群上的复函数都可以唯一展开为这些矩阵系数的线性组合。相比之下，不可约特征标是类函数空间 \(\operatorname{CF}(G)\) 的标准正交基；这个子空间的维数是共轭类数，通常小于 \(|G|\)。特征标只取各矩阵的对角元之和，全部矩阵系数则描述了更大的函数空间。

\ProseParagraphBreak
如果已经构造了一组两两不同构的不可约表示，且它们的次数平方和等于群阶，那么列表已经完备：任何遗漏的不可约还会给正则分解增加一个正的平方。反过来，平方和吻合本身不能证明候选表示不可约，也不能排除列表重复；这两项需要先核验。

\subsection{复群代数的矩阵块分解}
\label{b5:subsec:0203:02}

\begin{theorem}\label{b5:thm:complex-group-algebra}
设 \(G\) 为有限群，\(S_1,\ldots,S_r\) 为全部不同构的不可约有限维复表示，\(d_i=\dim_{\mathbb C}S_i\)，记 \(\rho_i\) 为 \(S_i\) 的作用及其在 \(\mathbb C[G]\) 上的线性延伸。由作用定义的复代数同态
\[
\Phi:\mathbb C[G]\longrightarrow\prod_{i=1}^r
\operatorname{End}_{\mathbb C}(S_i),\qquad
a\longmapsto\bigl(\rho_i(a)\bigr)_i
\]
为同构。因此 \(\mathbb C[G]\simeq\prod_i M_{d_i}(\mathbb C)\)。
\end{theorem}
\begin{proof}
各作用保持乘法和单位，故是代数同态。若 \(a\) 的全部像为零，它零化每个不可约，因而由正则分解也零化 \(\mathbb C[G]\)；作用于一得到 \(a=0\)。所以 \(\Phi\) 单射。两侧复维数分别为 \(|G|\) 和 \(\sum_i d_i^2\)，由\cref{b5:thm:regular-decomposition}相等，有限维线性代数使它满射。
\end{proof}

不选 \(S_i\) 的基时，右端写成自同态代数，作用同态本身有明确意义；写成具体矩阵代数则需要选基。一般底域下可能出现非平凡除代数，正特征整除群阶时还会失去半单性，所以这个复群代数公式不能原样用于所有域。

\ProseParagraphBreak
满射性还说明，各不可约类型上的作用可以彼此独立地指定：给定任意算子 \(A_i\in\operatorname{End}_{\mathbb C}(S_i)\)，都存在同一个群代数元素 \(a\)，同时满足 \(\rho_i(a)=A_i\)。所以不只是每个矩阵块单独得到全部线性算子，还可以让其余块为零，只在指定的一个块中实现所需的作用。

\ProseParagraphBreak
这里的 \(\operatorname{End}_{\mathbb C}(S_i)\) 与 Schur 引理中的 \(\operatorname{End}_G(S_i)\) 是不同的代数：前者容许所有线性算子，后者只容许与群作用交换的算子。例如，\(S_3\) 的二维标准表示 \(S\) 是三点置换空间中坐标和为零的子空间；它的特征标在三个共轭类上为 \(2,0,-1\)，范数平方为一，故不可约。于是 \(\operatorname{End}_G(S)=\mathbb C I\)，而刚证的同构使群代数在 \(S\) 上的像为全部 \(M_2(\mathbb C)\)。这个四维矩阵块作为左模也不是一个简单模：每个矩阵唯一分成只有第一列非零与只有第二列非零的两个矩阵，左乘保持这两部分，并在每一列上按 \(S\) 作用。因此它作为左 \(G\)-模为 \(S\oplus S\)，对应正则分解中的重数二。

\subsection{中心幂等元及块的半单情形}
\label{b5:subsec:0203:03}

\begin{theorem}\label{b5:thm:central-character-idempotents}
设 \(G\) 为有限群，\(S_1,\ldots,S_r\) 为全部不同构的不可约有限维复表示，其特征标为 \(\chi_i\)，\(d_i=\chi_i(1)\)。令
\begin{equation}\label{b5:eq:central-idempotent}
e_i=\frac{d_i}{|G|}\sum_{g\in G}\chi_i(g^{-1})g.
\end{equation}
则 \(e_i\in\mathbb C[G]\) 是两两正交的中心幂等元，\(\sum_i e_i=1\)，并且在 \(G\) 的任意有限维复表示 \(V\) 中，\(e_iV\) 正是类型 \(S_i\) 的等型分量。理想 \(\mathbb C[G]e_i\) 同构于 \(M_{d_i}(\mathbb C)\)。
\end{theorem}
\symidx{\(e_i\)：复群代数的本原中心幂等元}
\begin{proof}
为说明归一化系数的来源，先取
\[
z_i=\sum_{g\in G}\chi_i(g^{-1})g.
\]
因系数在共轭类上常值，\(z_i\) 位于中心。由 Schur 引理，它在不可约 \(S_j\) 上作用为某个标量 \(a_{ij}I\)，取迹并使用特征标正交性，得到
\[
d_j a_{ij}=\sum_{g\in G}\chi_i(g^{-1})\chi_j(g)
=|G|\langle\chi_i,\chi_j\rangle_G=|G|\delta_{ij}.
\]
因此 \(z_i\) 在 \(S_i\) 上作用为 \((|G|/d_i)I\)，在其他不可约表示上为零。要使前一个作用成为恒等算子，只需乘以 \(d_i/|G|\)，这就得到 \(e_i=(d_i/|G|)z_i\) 及\eqref{b5:eq:central-idempotent}。借\cref{b5:thm:complex-group-algebra}的单射性，作用上的平方、相互乘积与和的等式在群代数中也成立。对一般 \(V\) 用不可约分解，得到等型投影。最后 \(\Phi(e_i)\) 仅在第 \(i\) 个矩阵因子为单位，故相应理想恰是那个因子。
\end{proof}

不能再写成两个非零正交中心幂等元之和的非零中心幂等元，称为本原中心幂等元；由上面的矩阵乘积分解，各 \(e_i\) 恰为全部这类元素。它们对应的直和因子称为群代数的\term[qun-daishu-kuai]{块}[block of a group algebra]。在当前复半单情形，每个块只有一个不可约类型；模特征情形的块可能包含多个简单类型，不能把当前特点写成块的一般定义。

\ProseParagraphBreak
这里“本原中心”的限制不能省去。例如在 \(M_2(\mathbb C)\) 中，有 \(I_2=E_{11}+E_{22}\)，右边是两个非零正交幂等元，但它们不是中心元素。中心只有 \(\mathbb C I_2\)，其中的幂等元只有零和 \(I_2\)，所以 \(I_2\) 仍是本原中心幂等元。相应矩阵块不能再拆成非零的代数直积因子，却可作为左模分成两个简单左理想 \(M_2(\mathbb C)E_{11}\) 与 \(M_2(\mathbb C)E_{22}\)。这与刚才按列分解成两份二维表示一致：代数的中心块与左模的不可分解直和项是不同对象。

\ProseParagraphBreak
中心还可直接用共轭类和描述。对每个共轭类 \(C\)，令 \(z_C=\sum_{g\in C}g\)。一个群代数元素与全部群元素交换，当且仅当其系数在各共轭类上相同，所以这些 \(z_C\) 构成中心的一组基。另一方面，由矩阵乘积分解，每个中心元素在各块上是标量，因此 \(\{e_i\}\) 也是中心的一组基。取 \(c\in C\)，\(z_C\) 在 \(S_i\) 上为标量算子，其迹为 \(|C|\chi_i(c)\)，所以该标量为 \(|C|\chi_i(c)/d_i\)。比较各矩阵块上的作用，再用 \(\Phi\) 的单射性，得到两组中心基之间的换基公式
\[
z_C=\sum_i\frac{|C|\chi_i(c)}{d_i}\,e_i,
\qquad
e_i=\frac{d_i}{|G|}\sum_C\chi_i(c^{-1})z_C.
\]
第二式是将\eqref{b5:eq:central-idempotent}按共轭类合并。共轭类和按群元素组织中心，中心幂等元按不可约类型组织中心，特征标表给出两者的坐标转换；第五章将进一步考察类和乘法的整数系数。

\ProseParagraphBreak
例如在 \(\mathbb C[S_3]\) 中，令 \(T\) 为三个换位之和。这个共轭类大小为 \(3\)，换位在平凡、符号、二维标准表示 \(S\) 上的特征标值分别为 \(1,-1,0\)。代入上面的类和坐标公式 \(z_C=\sum_i |C|\chi_i(c)e_i/d_i\)，其三个系数分别为 \(3,-3,0\)，所以
\[
T=3e_{\rm triv}-3e_{\rm sgn}.
\]
再取 \(c=(123)\)，由 \(\chi_S=(2,0,-1)\) 得到
\[
e_S=\frac{2-c-c^2}{3}.
\]
取本原三次单位根 \(\omega\)。\(c\) 的作用满足 \(\rho_S(c)^3=I\)，且迹为 \(-1\)，故两个特征值为 \(\omega,\omega^2\)。在相应特征基下，群代数元素
\[
a=e_S\frac{c-\omega^2\cdot1}{\omega-\omega^2}
\]
在平凡、符号表示上均为零，在 \(S\) 上却作用为
\[
\rho_S(a)=
\begin{pmatrix}1&0\\0&0\end{pmatrix}.
\]
因而它确实只在二维矩阵块内取出一个坐标方向，具体实现了 \(\Phi\) 的满射性所保证的一种独立作用。
